% ============================================================
% III. METODOLOGÍA
% ============================================================

\section{Metodología}

El procedimiento experimental se dividió en tres grupos: animales, instrumentos musicales y voces humanas. Cada archivo de audio se almacenó en formato WAV y se procesó con el paquete \texttt{animal\_analyzer} del proyecto \texttt{WORK\_THREE}, apoyado en las funciones genéricas de \texttt{signal\_lab/spectral.py} y en NumPy, SciPy \cite{scipy} y Matplotlib.

Las tres grabaciones de animales (Gato, Perro, Gallo) se cargaron y analizaron de forma interactiva desde la interfaz gráfica \texttt{run\_animal\_gui.py} (pestañas dedicadas por categoría). Las funciones de análisis son independientes de la categoría, por lo que las grabaciones de instrumentos y voces se procesaron invocando directamente el mismo backend (\texttt{analyze\_audio}) de forma programática, sin pasar por esas pestañas específicas de la GUI.

El procedimiento general, tal como lo implementa \texttt{analyze\_audio}, fue:

\begin{enumerate}

\item Cargar el archivo de audio con \texttt{load\_audio}, que usa \texttt{librosa.load(path, sr=None, mono=True)} y preserva la frecuencia de muestreo nativa del archivo.

\item Registrar la frecuencia de muestreo y la duración total de la grabación.

\item Acotar la señal a los primeros $3$ s (\texttt{analysis\_window}, \texttt{ANALYSIS\_WINDOW\_S = 3.0}), dado que las grabaciones originales tienen una duración aproximada de $10$ s y solo esa ventana inicial se somete a análisis espectral y estadístico.

\item Representar la señal ventaneada en el dominio temporal.

\item Calcular la FFT de un solo lado con \texttt{compute\_fft} (\texttt{numpy.fft.rfft}, magnitud normalizada por $2/N$).

\item Obtener el espectro de magnitud y graficarlo.

\item Determinar la frecuencia dominante con \texttt{dominant\_frequency}, ignorando la componente en $0$ Hz.

\item Calcular la media y la desviación estándar con \texttt{signal\_stats}.

\item Calcular el puntaje de separabilidad estadística (Ec.~\ref{eq:separabilidad}) entre cada par de señales con \texttt{separability\_score} / \texttt{compare\_categories}.

\item Comparar los resultados obtenidos entre las diferentes señales.

\end{enumerate}

\subsection{Señales animales}

Se analizaron tres sonidos correspondientes a un gato, un perro y un gallo, cargados desde \texttt{audio\_samples/gato/}, \texttt{audio\_samples/perro/} y \texttt{audio\_samples/gallo/} respectivamente.

\begin{itemize}
\item Gato.wav
\item Perro.wav
\item Gallo.wav
\end{itemize}

% ============================================================
% FIGURAS ANIMALES
% ============================================================

\begin{figure}[H]
\centering
\includegraphics[width=\linewidth]{gato.png}
\caption{Señal temporal (primeros 3 s) y espectro frecuencial correspondiente al sonido de un gato.}
\label{fig:gato}
\end{figure}

\begin{figure}[H]
\centering
\includegraphics[width=\linewidth]{perro.png}
\caption{Señal temporal (primeros 3 s) y espectro frecuencial correspondiente al sonido de un perro.}
\label{fig:perro}
\end{figure}

\begin{figure}[H]
\centering
\includegraphics[width=\linewidth]{gallo.png}
\caption{Señal temporal (primeros 3 s) y espectro frecuencial correspondiente al sonido de un gallo.}
\label{fig:gallo}
\end{figure}

\begin{figure}[H]
\centering
\includegraphics[width=\linewidth]{animales_comparacion.png}
\caption{Comparación de los espectros frecuenciales de los sonidos animales.}
\label{fig:animales_comparacion}
\end{figure}

% ============================================================
% INSTRUMENTOS
% ============================================================

\subsection{Instrumentos musicales}

Para estudiar diferentes registros frecuenciales se seleccionaron tres instrumentos musicales:

\begin{itemize}
\item Contrabajo, como instrumento de baja frecuencia.
\item Piano, como instrumento de frecuencia media.
\item Flauta, como instrumento de frecuencia alta.
\end{itemize}

\begin{figure}[H]
\centering
\includegraphics[width=\linewidth]{contrabajo.png}
\caption{Señal temporal (primeros 3 s) y espectro frecuencial correspondiente al contrabajo.}
\label{fig:contrabajo}
\end{figure}

\begin{figure}[H]
\centering
\includegraphics[width=\linewidth]{piano.png}
\caption{Señal temporal (primeros 3 s) y espectro frecuencial correspondiente al piano.}
\label{fig:piano}
\end{figure}

\begin{figure}[H]
\centering
\includegraphics[width=\linewidth]{flauta.png}
\caption{Señal temporal (primeros 3 s) y espectro frecuencial correspondiente a la flauta.}
\label{fig:flauta}
\end{figure}

\begin{figure}[H]
\centering
\includegraphics[width=\linewidth]{instrumentos_comparacion.png}
\caption{Comparación de los espectros frecuenciales de los instrumentos musicales.}
\label{fig:instrumentos_comparacion}
\end{figure}

% ============================================================
% VOCES
% ============================================================

\subsection{Voces humanas}

Finalmente, se analizaron dos grabaciones correspondientes a dos personas diferentes. Ambas personas pronunciaron la misma frase con el propósito de controlar el contenido lingüístico de las señales.

La frase utilizada fue:

\begin{quote}
``Hoy estamos analizando señales de voz mediante la transformada de Fourier.''
\end{quote}

\begin{figure}[H]
\centering
\includegraphics[width=\linewidth]{persona1.png}
\caption{Señal temporal (primeros 3 s) y espectro frecuencial correspondiente a la Persona 1.}
\label{fig:persona1}
\end{figure}

\begin{figure}[H]
\centering
\includegraphics[width=\linewidth]{persona2.png}
\caption{Señal temporal (primeros 3 s) y espectro frecuencial correspondiente a la Persona 2.}
\label{fig:persona2}
\end{figure}

\begin{figure}[H]
\centering
\includegraphics[width=\linewidth]{personas_comparacion.png}
\caption{Comparación de los espectros frecuenciales de las dos voces humanas.}
\label{fig:personas_comparacion}
\end{figure}
